%%%%%%%%%%%%
%
% Installation and First Launch
%
%%%%%%%%%%%%

\chapter{Installation and First Launch}
\label{ch:installation-first-launch}

This chapter explains how to prepare the software environment and start the CPI forecasting dashboard for normal use. The focus is the user-facing launch workflow, not the internal implementation of the forecasting pipeline.

\section{Prerequisites}

Before starting the dashboard, the user should ensure that the following are available:

\begin{itemize}
    \item Python 3 with virtual-environment support,
    \item the repository folder including the \FILELOC{Code/} subdirectory,
    \item internet access for the initial dependency installation,
    \item a browser for opening the Streamlit dashboard.
\end{itemize}

\section{Folder Check Before Launch}

The most important working folder is \FILELOC{Code/}. Before installation or launch, verify that the following items exist:

\begin{enumerate}
    \item \FILELOC{Code/requirements.txt}
    \item \FILELOC{Code/run\_dashboard.bat} or \FILELOC{Code/run\_dashboard.sh}
    \item \FILELOC{Code/saved\_models/}
\end{enumerate}

If the saved model files are missing, the dashboard may still start but will not be able to display the expected forecasts. In that case, the issue should be treated as a maintainer-side problem and reported through the support process.

\section{Normal Installation Sequence}

\subsection{Windows}

For Windows users, the recommended first-run sequence is:

\begin{enumerate}
    \item open the \FILELOC{Code/} folder,
    \item run \SHELL{setup\_env.bat},
    \item after installation completes, run \SHELL{run\_dashboard.bat}.
\end{enumerate}

\subsection{macOS / Linux}

For macOS or Linux users, the recommended first-run sequence is:

\begin{enumerate}
    \item open a terminal in \FILELOC{Code/},
    \item run \SHELL{./setup\_env.sh},
    \item after installation completes, run \SHELL{./run\_dashboard.sh}.
\end{enumerate}

\section{Alternative Start Command}

The repository also provides \FILELOC{start.bat} and \FILELOC{start.sh}. These scripts are useful when the full environment and dashboard startup should be executed in one sequence. However, for normal day-to-day usage, the dedicated dashboard start scripts are easier to understand and are therefore preferred in this manual.

\section{Expected First Launch Result}

After a successful launch, Streamlit prints a local address in the terminal. Open that address in a browser. The first screen should show:

\begin{itemize}
    \item the dashboard title,
    \item the configuration sidebar,
    \item forecast metrics,
    \item a forecast graph,
    \item explanatory sections below the main chart.
\end{itemize}

\begin{tcolorbox}[colback=orange!5!white, colframe=orange!75!black, title=Important Usage Note]
Normal users should \textbf{not} update the raw CPI file, retrain the forecasting models, or edit Python modules just to use the dashboard. If such actions become necessary, the issue belongs to a maintainer and should be handled through the support process described later in this manual.
\end{tcolorbox}
